\section{Introduction}
\label{sec:intro}

% Motivating example
Consider a three-stage AI agent pipeline: a transcription agent converts audio to text,
a summarization agent condenses the transcript, and a report-generation agent produces
the final output. Each agent is compensated via a continuous payment stream proportional
to its throughput. When the summarization agent reaches its GPU capacity limit, the upstream
payment stream continues---but the surplus cannot be ``dropped'' like excess data packets.
Money must go somewhere.

% Problem statement
This scenario illustrates a fundamental gap in agent-to-agent
payment protocols. Recent systems---Google's AP2, Coinbase's x402~\cite{coinbase2025x402},
OpenAI--Stripe's ACP~\cite{openai2025acp}, and Visa's TAP~\cite{visa2025tap}---all address
\emph{authorization} and \emph{trust} for agent payments, but none provide \emph{flow control}:
the ability to dynamically route, throttle, or redistribute monetary flows based on real-time
capacity constraints. Traditional payment systems assume discrete, settled transactions; streaming
payment protocols like Superfluid~\cite{superfluid2024} enable continuous flows but have no
built-in congestion awareness.

% Gap
In data networks, this problem was solved decades ago. Backpressure routing, introduced by
Tassiulas and Ephremides~\cite{tassiulas1992stability}, provides throughput-optimal scheduling
by using queue differentials to guide packet routing. Kelly's proportional fairness
framework~\cite{kelly1998rate} established that shadow prices on capacity constraints can be
interpreted as economic signals. Yet no prior work has unified these insights with monetary
mechanism design for agent economies.

% Contributions
We introduce \emph{Backpressure Economics} (BPE), a cryptoeconomic mechanism that adapts
backpressure routing from communication networks to monetary flows. Our contributions are:

\begin{enumerate}
    \item \textbf{Formal model} (\Cref{sec:model}): We define a network payment flow model
    where sinks signal capacity via EWMA-smoothed declarations, and payments are routed via
    a max-weight scheduling rule adapted from Tassiulas--Ephremides.

    \item \textbf{Throughput optimality} (\Cref{sec:throughput}): We prove that BPE achieves
    throughput-optimal payment allocation within the capacity region, using a Lyapunov drift
    argument modified to handle the no-drop monetary constraint via a bounded overflow buffer.

    \item \textbf{Protocol design} (\Cref{sec:protocol}): We implement BPE using Superfluid's
    General Distribution Agreement (GDA) on Base, with commit-reveal capacity signaling,
    stake-weighted Sybil resistance, and multi-stage pipeline composition.

    \item \textbf{Simulation} (\Cref{sec:evaluation}): We evaluate BPE in a 50-sink,
    10-source agent economy, demonstrating 95.7\% allocation efficiency versus 93.5\% for
    round-robin and 79.7\% for random allocation.
\end{enumerate}
